\section{Introduction}
\label{sec:introduction}

Behavior-Driven Development (BDD) asks teams to capture expected system behavior as structured, human-readable scenarios, typically written in Gherkin and tied back to agile user stories~\cite{cucumber2023gherkin}.
Authoring these scenarios by hand is slow going: someone has to take an informal requirement and turn it into a syntactically valid Given--When--Then sequence, all while keeping the domain vocabulary and testable intent intact.
Rathnayake et al.~\cite{rathnayake2026llm} recently released 500 paired user stories and manually authored Gherkin scenarios drawn from IntelligenceBank products --- a dataset that underscores both how useful good reference scenarios are and how much effort goes into producing them at scale.

Large language models (LLMs) have already shown some promise for automating natural-language software engineering artifacts, from requirements refinement to test-related text generation.
On the Rathnayake et al. dataset, GPT-4 evaluated zero-shot earned a mean expert rating of 4.63 out of 5.0~\cite{rathnayake2026llm}, which hints that current models can generate plausible Gherkin from user stories with no task-specific training at all.
Expert ratings are only part of the picture, though; they say little about whether a model's output clears the kind of fixed, reproducible quality gates a team would actually want to embed in a continuous evaluation pipeline.

Few prior studies on LLM-based BDD generation bring together \emph{automated} quality metrics, \emph{statistical hypothesis testing}, and \emph{fixed acceptance thresholds} within one empirical design --- in our review of 13 closely related studies, seven were missing at least one of these three elements, which makes cross-study comparison and clean pass/fail judgments hard to come by.
As a result, it's still an open question how far zero-shot generation can be trusted as a stand-in for manual scenario authoring once explicit quality contracts are on the table.

We take up that question here, through a controlled, replication-style evaluation of GPT-4o on a curated subset of the Rathnayake et al. dataset.
This work contributes the following:
\begin{itemize}
  \item An evaluation --- the first we are aware of --- of \textbf{GPT-4o zero-shot Gherkin generation} from user stories using automated metrics: cosine semantic similarity ($\geq 0.85$) via \texttt{all-MiniLM-L6-v2} embeddings~\cite{reimers2019sentence}, and executable syntax rate ($\geq 80\%$), each assessed with one-tailed statistical tests against pre-registered thresholds.
  \item Results from $N{=}100$ pairs across four IntelligenceBank domains, where \textbf{RQ1 (semantic similarity) fell short of its threshold} while \textbf{RQ2 (executable syntax) met it}.
  \item A \textbf{public replication package} --- sampled data, generation scripts, API logs, analysis notebooks, and statistical summaries --- available at \url{https://github.com/khanhkhanh55355/RT-SWT-003-nhom3}.
\end{itemize}

The remainder of the paper unfolds as follows: Section~\ref{sec:related-work} surveys related work on LLM-based BDD and Gherkin generation; Section~\ref{sec:method} lays out our research questions and evaluation metrics; Section~\ref{sec:experimental-setup} covers the dataset, model configuration, and analysis procedure; Section~\ref{sec:results-discussion} presents and discusses the results; Section~\ref{sec:threats-validity} addresses threats to validity and how we mitigated them; and Section~\ref{sec:conclusion} closes with conclusions and future directions.
